% Shared macros for the lab book. Included by labbook.tex before \begin{document}.

% ---------------------------------------------------------------------------
% Evidence tags. These are the LaTeX counterparts of the markdown tags
% [S:...] [D:...] [C:...] [I] [?] defined in PROVENANCE.md. Every factual
% sentence in a shard carries one. Untagged means unverified.
% ---------------------------------------------------------------------------

\definecolor{evsourced}{RGB}{0,102,0}
\definecolor{evderived}{RGB}{0,0,153}
\definecolor{evcomputed}{RGB}{0,102,153}
\definecolor{evinferred}{RGB}{153,76,0}
\definecolor{evopen}{RGB}{178,0,0}

% \evS{citekey}{location} -- stated by a registered source at that location.
% Also emits \cite so the source appears in the bibliography. Both arguments
% may be empty, which renders a bare [S] for referring to the tag in prose.
\newcommand{\evS}[2]{%
  \textcolor{evsourced}{\footnotesize[S%
    \ifx\relax#1\relax\else:\,\cite{#1}\fi
    \ifx\relax#2\relax\else\,#2\fi]}}

% \evD{derivation-id} -- our derivation, checked, under derivations/.
\newcommand{\evD}[1]{\textcolor{evderived}{\footnotesize[D:\,\texttt{#1}]}}

% \evC{run-id} -- our computation, recorded under runs/.
\newcommand{\evC}[1]{\textcolor{evcomputed}{\footnotesize[C:\,\texttt{#1}]}}

% \evI -- our inference or interpretation; plausible, not yet checked.
\newcommand{\evI}{\textcolor{evinferred}{\footnotesize[I]}}

% \evQ -- needs a source or a check.
\newcommand{\evQ}{\textcolor{evopen}{\footnotesize[?]}}

% ---------------------------------------------------------------------------
% Status and gates
% ---------------------------------------------------------------------------

% \gates{S,D,M} -- which validation gates were actually run for a result.
\newcommand{\gates}[1]{\textsf{\footnotesize gates:\,#1}}

% \status{active} -- status of a section's content.
\newcommand{\status}[1]{\textsf{\footnotesize status:\,#1}}

% \TODO{...} -- visible placeholder; never leave one in a promoted result.
\newcommand{\TODO}[1]{\textcolor{evopen}{\textbf{TODO:} #1}}

% \pathid{runs/2026-07-29-...} -- typewriter text that may break at its own
% punctuation. Run ids and DOIs are wider than the table columns that hold
% them, and plain \texttt has no break points, so they crossed the margin.
% Use this, not \texttt, for any path, run id or DOI inside a table cell.
% \DeclareUrlCommand, not a \newcommand wrapping \Url: the argument has to be
% read with url's own catcodes, which a normal macro argument has already lost.
\usepackage{url}
\DeclareUrlCommand\pathid{\urlstyle{tt}}

% ---------------------------------------------------------------------------
% Physics shorthand. Add sparingly, and only for notation fixed in
% CONVENTIONS.md.
% ---------------------------------------------------------------------------

\newcommand{\ket}[1]{\lvert #1 \rangle}
\newcommand{\bra}[1]{\langle #1 \rvert}
\newcommand{\ev}[1]{\langle #1 \rangle}
\newcommand{\braket}[2]{\langle #1 \vert #2 \rangle}
\newcommand{\dd}{\mathrm{d}}
